\documentclass[sigconf,nonacm,balance=false]{acmart}

\settopmatter{printacmref=false}
\setcopyright{none}
\acmConference[AI 231 ME3]{AI 231 ME3: Voice Duplex}{October 2026}{Mapua University}
\acmBooktitle{AI 231 ME3: Voice Duplex}
\acmDOI{}
\acmISBN{}

\usepackage{booktabs}
\usepackage{tabularx}
\usepackage{listings}
\usepackage{xcolor}
\usepackage{enumitem}
\usepackage{microtype}

\definecolor{codebg}{RGB}{245,247,250}
\definecolor{codecomment}{RGB}{75,90,105}
\lstset{
  basicstyle=\ttfamily\scriptsize,
  backgroundcolor=\color{codebg},
  frame=single,
  rulecolor=\color{black!12},
  breaklines=true,
  columns=fullflexible,
  commentstyle=\color{codecomment},
  showstringspaces=false
}

\title{ME3 Voice Duplex: A Tool-Using, Full-Duplex Voice Agent for Raspberry Pi}

\author{Crizepvill F. Dumalaog}
\affiliation{\institution{Mapua University}\department{Master of Engineering in Artificial Intelligence}\city{Manila}\country{Philippines}}
\email{202521406}

\begin{document}

\begin{abstract}
ME3 is a Raspberry Pi voice-agent prototype that combines local speech processing,
full-duplex audio, a provider-selectable remote language-model brain (Qwen by
default, with optional GPT-5.6 Luna), and project-owned Model
Context Protocol (MCP) services. The system is designed for an important
interaction beyond question answering: a user can ask by voice for a web
application or other coding task, and the model can call bounded file, search,
command, and background-job tools to perform the work. No task-specific model
training is required. Instead, the system integrates pretrained Silero VAD,
Qwen2-Audio-7B-Instruct ASR, Qwen3-TTS-12Hz-1.7B-CustomVoice TTS, and a
retained but inactive ElevenLabs adapter, together with an OpenAI-compatible
Qwen/vLLM endpoint or the OpenAI GPT-5.6 Luna API. The active Qwen audio path
uses protected vLLM listeners; the text ToolBrain and bounded tools remain
shared within the final release.
This paper describes the architecture, tool-call protocol, safety boundary,
local verification, Raspberry Pi release workflow, and evidence boundary for
the separate FULL-DUPLEX-BENCH v1.5 evaluation. The current implementation is
an engineering release, not a claim of completed paper-grade overlap results.
The current revision also specifies the Vocalis-Pi event contract: continuous
VAD/transcript inputs, explicit turn states, interrupt controls, and an
adapter seam for asynchronous MCP harness work. The workstation acceptance
includes the existing local checks plus mocked request-shape and configuration
tests for the Qwen audio profile. The protected DGX3 listeners were
authenticated and exercised over the restored VPN route. Qwen2-Audio ASR and
Qwen3-TTS both returned provider-live audio results, and a complete injected
voice/tool turn passed. The latest accepted Qwen release was deployed to
\texttt{/home/dalmacio/Desktop/dandan-qwen}; the former Own deployment was
removed from the active Pi path and moved to a dated archive. The live Qwen
service returned HTTP 200, loaded its Qwen audio models, and reported a
nonzero Fifine startup signal. A clean human spoken turn and paper-native
overlap run remain separate evidence gates. The vendored benchmark adapter,
LAN browser-speaker server, browser music handoff, and operator restart
control are part of the final Qwen package.
\end{abstract}

\keywords{voice agents, full duplex, Raspberry Pi, tool calling, MCP, Qwen,
GPT-5.6 Luna,
speech interaction, edge computing}

\maketitle

\section{Introduction}

Voice interfaces are often presented as a sequence: wait for the user to stop,
transcribe one utterance, generate text, and play a response. A useful
assistant must also manage interruption and action. A user should be able to
say ``build a web app'' and receive evidence of files and verification, rather
than a plausible paragraph describing code that was never created.

ME3 addresses these requirements with a local audio loop and a remote tool-using
language model. The design separates four concerns. First, the Pi keeps
voice-sensitive processing local through VAD, ASR, and TTS. Second, a Qwen
model served by an OpenAI-compatible vLLM API performs dialogue reasoning by
default; GPT-5.6 Luna is an optional OpenAI-compatible alternative.
Third, typed tools provide an auditable action surface. Fourth, a release and
deployment workflow distinguishes local preparation from actual Pi installation.

The project is intentionally not a new model-training project. The agent
behavior comes from inference and orchestration: pretrained speech models,
Qwen or GPT-5.6 Luna inference, MCP/tool dispatch, and a bounded local execution policy. This
choice reduces iteration time and makes the engineering questions explicit.

The contributions of this work are:
\begin{itemize}[leftmargin=*]
\item a full-duplex local voice loop with barge-in during TTS playback;
\item a provider-neutral OpenAI-compatible tool loop that can create and verify
a static web application through voice or text;
\item one bounded registry shared by direct Qwen calls and MCP wrappers;
\item a confirmation-gated Telegram Bot API text adapter with a safe
provider-response boundary, while ordinary phone calls and SMS remain out of
scope;
\item background task supervision so builds do not block the microphone;
\item a reproducible local-dev, package, deploy, and evidence workflow for the
Raspberry Pi target, including a LAN browser-speaker transport and a vendored
Full-Duplex-Bench v1/v1.5 adapter.
\end{itemize}

\section{System Requirements and Scope}

The required runtime contains four owned MCP services: VAD, ASR, TTS, and
tasks. Pipecat is retained as the harness-facing integration path, while the
local loop is the reference implementation used for Windows development and
the Pi service. The runtime configuration is kept in one YAML file so the
audio rate, speech thresholds, model choices, endpoint, tool budget, and
deployment path are not duplicated across scripts.

The current engineering scope includes interruptible spoken conversation,
remote Qwen inference, project-local coding tools, public web lookup, timers,
system information, the existing ME2 bridge, and credential-gated Telegram
text and optional email actions. It does not claim unrestricted
Onit equivalence. An always-listening device needs a narrower default action
surface than an interactive coding IDE. Capabilities that would add external
side effects, credential use, repository rewrites, or arbitrary shell access
must be added as reviewed typed tools.

\section{Architecture}

\subsection{Audio and conversation path}

The default audio path is:
\begin{center}
\texttt{microphone -> Silero VAD -> ASR -> Qwen/Luna -> fixed profile TTS -> speaker}.
\end{center}
Capture remains active while TTS plays. The local state machine has four states:
\texttt{idle}, \texttt{listening}, \texttt{thinking}, and \texttt{speaking}.
Speech during \texttt{speaking} stops playback and starts a new turn. A short
pre-roll is retained so the beginning of an utterance is not clipped. The VAD
endpoint requires a minimum speech interval followed by configured trailing
silence.

The active Qwen profile always uses Qwen3-TTS CustomVoice \texttt{aiden} with a formal style
instruction. Replies, backchannels, and
completion announcements use the same profile voice. The Philippines-market
profile keeps English, Filipino, and Taglish in the supported interaction set.
The spoken policy is formal and professional: the fixed long-task acknowledgement
is ``I am processing your request now,'' and the assistant is instructed to
avoid flirtatious, intimate, suggestive, teasing, or overly familiar language.
The ElevenLabs adapter is retained as a documented future experiment but is
not auto-selected by the active demo runtime.
Qwen remains the default reasoning and tool-calling model; GPT-5.6 Luna is an
optional provider using the same tool loop. ElevenLabs Realtime v1
Mini/Max is a separate hosted voice-agent architecture that also manages the
LLM and ASR settings, so it is not silently substituted into this experiment.

On 26 October 2026, an authorized local provider-only test used the
gitignored Windows development credential and an operator-selected account
voice. ElevenLabs returned 54 raw-PCM chunks (163,468 bytes at 16~kHz), and a
Filipino-language test phrase played through the workstation's default
speaker. The adapter was corrected to reassemble 16-bit samples split across
HTTP chunk boundaries and to classify chunked-encoding failures using the
correct Requests exception handling. The account voices inspected
for this test were English-labelled; therefore this result verifies the
provider integration and output path, not approval of a native Filipino
accent.

The duplex state machine is independent of the Qwen model. This makes the
interaction control testable without a network request and leaves room for a
Pipecat transport to reuse the same transition logic.

For the ME2-style Pi web experience, capture remains on the Pi while an
optional browser-output transport publishes generated 16~kHz mono PCM over a
LAN WebSocket. The browser schedules that stream through Web Audio and can
select the PC's speaker; alternatively, the same UI selects a Pi-local output
device. The HTTP dashboard exposes safe PortAudio device names and restarts
only the audio loop when a route changes. A one-click browser connection is
required by autoplay policy, but it does not create a push-to-talk interaction:
the Pi microphone remains continuously active and barge-in is handled by the
Pi state machine.

The selectable \texttt{qwen\_instruct} profile changes only the ASR/TTS
providers. Local Silero VAD and the same duplex state machine accept barge-in.
Completed turns are sent to \texttt{Qwen/Qwen2-Audio-7B-Instruct} through an
OpenAI-compatible multimodal \texttt{/v1/chat/completions} request carrying
the WAV as an \texttt{audio\_url} data URL. Replies are rendered by
\texttt{Qwen/Qwen3-TTS-12Hz-1.7B-CustomVoice} through
\texttt{/v1/audio/speech}. The text ToolBrain remains the existing
\texttt{Qwen/Qwen3.8-27B} or optional Luna provider, so requests such as
``build a web app'' retain the same file and verification tools. Endpoint and
key values are runtime-only. An ASR/TTS pair is therefore not a replacement
for the general text tool-calling model.

\subsection{Vocalis-Pi event contract}

The implementation makes the full-duplex behavior observable through a
transport-independent Vocalis-Pi orchestrator. Its external states are
\texttt{LISTENING}, \texttt{THINKING}, \texttt{SPEAKING},
\texttt{BACKCHANNELING}, and \texttt{INTERRUPTED}. The local loop remains
active on the microphone while the speaker plays. A VAD update received during
\texttt{SPEAKING}, \texttt{THINKING}, or \texttt{BACKCHANNELING} publishes a
\texttt{barge\_in\_event}, then emits \texttt{FLUSH\_QUEUE} and
\texttt{FADE\_OUT} TTS controls before returning to \texttt{LISTENING}.

The protocol carries four input/event types: VAD state, transcript deltas,
overlap signals, and asynchronous harness events. Its output types are
assistant-text deltas, TTS control actions, turn-state transitions, and
optional harness dispatch requests. A long Qwen tool turn can therefore
produce a short \texttt{BACKCHANNELING} acknowledgement while its task
continues. Tool progress is mirrored as sanitized harness events; arguments
and results are excluded from these events because they can contain user data
or credentials. The exact channel names and JSON examples are documented in
the repository's Vocalis protocol manual.

The optional harness adapter emits a JSON-RPC 2.0 \texttt{tools/call} request
named \texttt{harness\_dispatch}. The default verified release keeps the
bounded Qwen tool loop in-process for low overhead and uses the adapter as a
stable seam for an Onit, Pipecat, or OpenCode background worker. This makes
the contract testable without claiming that a separate external harness is
already running in every deployment.

\subsection{Provider-neutral tool-using brain}

The remote brain is an OpenAI-compatible chat-completions client. The first
request contains the recent dialogue and JSON Schema definitions for the ME3
tools. Qwen or GPT-5.6 Luna may return assistant text or a list of function
calls. For every
call, ME3 parses the arguments, executes the local function, appends a
\texttt{role=tool} message, and asks Qwen to continue. The checked-in
configuration allows 24 tool rounds for demanding multi-file or web-app tasks.
After the final permitted batch, ME3 makes one no-tool summarization request
instead of exposing an internal tool-limit error to the user.

For Qwen, the checked-in model is served by vLLM. For Luna, the model ID is
\texttt{gpt-5.6-luna}; hidden reasoning is disabled, temperature remains
omitted, and the client requests a direct response stream when supported. The
protocol is illustrated below.

\begin{lstlisting}
{
  "model": "Qwen/Qwen3.8-27B",
  "messages": [{"role":"user", "content":"build a web app"}],
  "tools": [{"type":"function", "function":
    {"name":"write_file", "parameters": {}}}],
  "tool_choice": "auto"
}
\end{lstlisting}

After a tool call, the result is returned as structured JSON. Qwen therefore
receives the actual file path, byte count, return code, or error. The final
spoken response is not accepted as evidence by itself.

\subsection{Owned MCP services}

The VAD server exposes frame and stream detection plus endpointing. The ASR
server exposes PCM and file transcription. The TTS server exposes synthesis,
streaming chunks, and installed voices. The tasks server retains timers,
system information, and the ME2 bridge, and now wraps the same bounded project
tool registry used by the direct Qwen path. The shared implementation avoids
one policy for the local loop and another policy for MCP clients.

\section{Tool Registry and Safety Boundary}

Table~\ref{tab:tools} summarizes the direct functions visible to Qwen.

\begin{table}[t]
\caption{Direct Qwen tools in the ME3 registry}
\label{tab:tools}
\small
\begin{tabularx}{\columnwidth}{@{}lX@{}}
\toprule
Tool & Purpose \\
\midrule
\texttt{list\_files} & List project files and sizes. \\
\texttt{read\_file} & Read bounded project text. \\
\texttt{write\_file} & Create or replace a project text file. \\
\texttt{search\_text} & Search project text files. \\
\texttt{run\_command} & Run a direct allowlisted command. \\
\texttt{start\_background\_task} & Start a bounded build or test. \\
\texttt{task\_status/output/cancel} & Supervise a background task. \\
\texttt{system\_info} & Return non-secret host information. \\
\texttt{fetch\_url/web\_search} & Read public HTTP(S) resources. \\
\texttt{git\_status} & Inspect repository status without mutation. \\
\texttt{channel\_health} & Report safe Telegram and optional email configuration. \\
\texttt{send\_telegram\_message/send\_email} & Confirmation-gated external
communication through Telegram Bot API and optional Resend. \\
\bottomrule
\end{tabularx}
\end{table}

The path policy resolves every path inside the ME3 project. Credential-like
names, environment files, private keys, the Git directory, the virtual
environment, model caches, and generated release trees are protected. Reads
and writes are size bounded. Commands are tokenized and run without a shell;
shell operators and patterns associated with deletion, privilege escalation,
SSH/SCP, repository reset/clean/push, and package publication are rejected.
Long commands run with standard input closed, captured output, and a timeout.

This boundary is a design choice rather than an implementation accident. A
voice transcript can be ambiguous, and the assistant is not an appropriate
place for silent arbitrary shell execution. The registry is sufficient for the
demonstration web app while keeping higher-risk operations explicit.

The external communication tools use the same registry and are stricter than
the local project tools. \texttt{channel\_health} is network-free. Each
\texttt{send\_telegram\_message} or optional \texttt{send\_email} request first
validates its target and returns a confirmation-required status unless the
operator has explicitly confirmed the side effect. The adapters then send one
bounded REST request and return only a status, channel, and safe provider ID;
provider response bodies and authorization headers are discarded. Ordinary
Telegram voice calls are not available through the Bot API. Inbound webhooks,
identity/CRM actions, consent, recording, and human transfer remain outside
this adapter boundary.

\section{Web-App Capability}

The target workflow is a dependency-free web application when the user does not
name a framework. Qwen writes an application directory under
\texttt{agent\_workspace/<name>}, lists the result, reads the important files,
and runs a bounded verification command if one is appropriate. The provided
demonstration contains \texttt{index.html}, \texttt{styles.css}, and
\texttt{app.js}; it provides a command input, microphone placeholder, status
area, and tool-activity log.

The capability is demonstrated twice: the first live request used Qwen to
create the HTML file through a real tool call; a later live request used Qwen
to list the directory and read all three files. This demonstrates inference,
tool dispatch, file mutation, and evidence retrieval without training a model.
The demo is intentionally local and static; the current safety policy does not
publish it to the public internet or install an arbitrary server.

\section{Deployment Workflow}

ME3 follows a four-stage local-dev workflow:
\begin{enumerate}[leftmargin=*]
\item \textbf{Environment gate:} import the shared Python dependencies.
\item \textbf{Smoke gate:} load all four MCP services, call their health tools,
construct the engineering benchmark, and query PortAudio when available.
The vendored public fd-benchmark is also available through a direct-feed ME3
adapter; its simulation is used as a protocol smoke test, not as a hardware
result.
\item \textbf{Package:} copy the runtime files into a versioned release, write
SHA-256 entries, and create a zip archive.
\item \textbf{Deploy:} transfer the archive to the configured Pi, merge its
stable \texttt{me3\_voice\_duplex} directory, run the idempotent installer,
and record a deployment receipt.
\end{enumerate}

\begin{sloppypar}
The current target is \url{dalmacio@cfdfnjrpi5.local} at
\url{/home/dalmacio/Desktop/dandan}; the Qwen profile is deployed separately
at \url{/home/dalmacio/Desktop/dandan-qwen}. The separate ME2 deployment is
kept outside the ME3 release. Both current receipts report \texttt{deployed}
after SSH, SCP, extraction, and installation. The active Qwen service returns
HTTP 200, its Qwen audio models load, and the Fifine capture path reports a
nonzero startup signal. This is Pi service/audio-start evidence, not yet a
clean human spoken-turn or speaker-observation claim.
\end{sloppypar}

The API key is never included in the archive. On the Pi it must be provisioned
through a secure runtime environment. The Windows DPAPI helper can report
configured/not-configured and can provide an in-memory child-process value for
local testing, but its secret is not printed or written to the repository.

\section{Verification and Evidence}

The current local verification sequence is:

\begin{lstlisting}
python scripts/dev_env_check.py
python scripts/local_smoke_test.py --with-audio
python scripts/check_paper_alignment.py
python -m pytest tests -q
python scripts/run_me3_module.py me3.agent --brain remote --text \\
  "Use system_info and summarize the host."
\end{lstlisting}

The active Qwen test suite passed 84 tests in the current checkout. The
tool-focused tests cover OpenAI-compatible schemas, blocked shell patterns,
background task completion, and a mocked multi-step Qwen response. The
ElevenLabs tests additionally cover split-byte PCM reassembly, genuinely odd
total PCM rejection, timeout retry, and chunked-encoding retry without
credential disclosure. The channel tests cover confirmation gating, Telegram
chat-id validation, Bot API request construction, Resend request construction,
and safe metadata. Python compilation also passed.
The Vocalis tests cover state transitions, JSON-RPC harness dispatch,
barge-in TTS-control ordering, and LocalVoiceLoop tool-progress bridging.
One current injection-only fixture run also completed with one turn and zero
errors; it remains engineering evidence and not physical-speaker or v1.5
evidence.
The live Qwen smoke test records an actual tool call and final answer. The
local smoke gate loads the four MCP servers and confirms the current Windows
PortAudio inventory. The integrated fd-benchmark simulation also completed a
10-trial development-dataset run, yielding a protocol-and-report smoke result
with no ASR judge; this is not a paper-grade or physical Pi result. These are
engineering checks, not a claim about Pi audio.

The evidence status at this revision is summarized in
Table~\ref{tab:progress}.

\begin{table}[t]
\caption{ME3 progress evidence at the 6 November 2026 revision}
\label{tab:progress}
\small
\begin{tabularx}{\columnwidth}{@{}lX@{}}
\toprule
Evidence lane & Status and interpretation \\
\midrule
Tool-capable Qwen & Verified in earlier live runs: Qwen created and inspected
a web-app artifact through bounded tools. \\
Local duplex audio & Verified with the Windows microphone, Piper output, and
two barge-in cancellations; echo can contaminate ASR on loud speakers. \\
LAN browser audio & The Pi-side HTTP/WebSocket server, browser PCM sink, device
selection endpoint, and ME2-style SSH forwarding pass local protocol tests;
physical Pi capture and LAN audio are still separate hardware gates. \\
Integrated fd-benchmark & Vendored MIT benchmark source and ME3 direct-feed
adapter pass import/protocol checks; a 10-trial development-dataset simulation
completed without ASR/judge scoring. This is not a paper-grade result. \\
Active Qwen Instruct audio profile & Adapter, runtime selection, configuration checks,
and request-shape tests pass. Over the VPN, Qwen2-Audio ASR returned in 431.9 ms
and Qwen3-TTS returned in 851.6 ms; the same matched run recorded
faster-whisper ASR at 2732.0 ms and Piper TTS at 100.1 ms TTFB/176.7 ms total.
The injected Qwen voice/tool lane completed one turn with zero errors and five
shared tool identities. \\
ElevenLabs renderer & Verified locally: provider-live 16-kHz PCM playback
through the default speaker; not a native-Filipino voice-quality result. \\
Qwen endpoint & Authenticated model discovery and a streamed text probe pass
from the deployed Pi; the operator has also confirmed a human Pi Fifine turn
with a browser-speaker reply. Automatic physical-speaker barge-in remains a
separate gate. \\
External channels & Telegram/Resend adapters and fake-request tests pass; no
live Telegram message or email side effect has been placed, and phone calls/SMS
are out of scope. \\
Pi hardware & The latest Pi process loaded the Fifine capture path and reported
a nonzero startup signal; the operator-confirmed human spoken turn is complete;
automatic physical barge-in and echo behavior remain open. \\
Paper-native FDB v1.5 & Incomplete: no clean 499-utterance overlap artifact
or paper-grade result is claimed. \\
\bottomrule
\end{tabularx}
\end{table}

The deployed Pi preflight is also an explicit hardware gate. It is read-only,
does not alter mixer state, and returns exit code 2 when no capture device is
visible. This keeps a playback-only Pi from being mislabeled as a working
full-duplex voice endpoint.

When speaker hardware is unavailable, the project uses a separate
injection-only lane. The evaluator feeds recorded 16~kHz PCM into the same
VAD, ASR, Qwen/tool, TTS, and interruption state machine at wall-clock frame
cadence, while a virtual sink records TTS timing without opening PortAudio
output. This verifies injected-input agent behavior and virtual barge-in, but
does not claim physical-speaker behavior or paper-native v1.5 evidence.

\begin{sloppypar}
To align the class evaluation with the supplied public benchmark, the release
vendors the MIT-licensed fd-benchmark repository and adds
\texttt{scripts/me3\_fdbench\_agent.py}. The adapter connects the real
\texttt{LocalVoiceLoop} to the benchmark's direct-feed WebSocket: incoming
16~kHz mono int16 frames enter the same VAD/ASR/tool/TTS path, outgoing PCM is
recorded by the benchmark rig, and a barge-in emits its queue-clear event. A
small simulation run verifies the integration without opening a microphone or
speaker; full datasets, ASR/judge scoring, and a clean Pi run remain required
for publication-quality measurements.
\end{sloppypar}

The project also contains a separate engineering latency driver and a
paper-native FULL-DUPLEX-BENCH v1.5 contract. The v1.5 contract includes four
scenarios with 200, 99, 100, and 100 stimuli and response categories RESPOND,
RESUME, UNCERTAIN, and UNKNOWN. The engineering five-scenario driver must not
be presented as the paper's 499-utterance result. A complete captured overlap
artifact, transcript classification, and prosodic before/after analysis are
still required before claiming a paper-grade result.

\section{Discussion}

The main engineering trade-off is locality versus model capability. Keeping
speech I/O local reduces the need to upload audio and makes the Pi responsive,
while Qwen supplies stronger reasoning and code generation than a small local
model. The cost is network dependence, endpoint latency, and the need to
provision the API key on the target device. The typed registry makes the
boundary visible and testable, but it is not a complete replacement for an
interactive coding environment.

The local ElevenLabs result narrows one implementation risk: the selected
streaming endpoint can produce audio in the PCM format consumed by the
full-duplex playback path, and the adapter tolerates arbitrary HTTP chunk
boundaries. It is not part of the active single-voice runtime. A complete
voice-agent turn still requires both a reachable Qwen brain and an operator
approved voice profile for the intended Filipino customer-support deployment. The new
communication adapters narrow the integration boundary but do not establish
that the agent can safely replace a staffed call center: provider accounts,
webhooks, media streaming, consent, retention, escalation, and supervised
quality gates remain necessary. Provider rendering, Qwen/tool execution, and
external-channel side effects are therefore reported as separate evidence
lanes rather than collapsed into one end-to-end claim.

The full-duplex behavior is also more than a transport setting. VAD must keep
running during playback, the audio output must be cancellable, and the
turn-state transition must preserve the interrupted user's utterance. Echo
control, microphone placement, and ALSA device selection remain physical Pi
concerns. A successful Windows loopback does not prove those conditions.

Future work should add a Pi-local live self-test, Telegram inbound webhook
signature verification, secure credential provisioning and rotation, an
operator-owned Telegram/Resend test, a reviewed web-app preview server with
explicit approval, automatic physical-speaker barge-in/echo evaluation, and the full v1.5
overlap capture/evaluation. Additional Onit-like capabilities should be
introduced as typed, auditable functions rather than by removing the voice
safety boundary.

\section{Conclusion}

\noindent ME3 now combines full-duplex voice interaction with a real Qwen tool loop and
an explicit Vocalis-Pi event contract. A
spoken build request can result in actual project files and verification, and
the agent can supervise longer work without freezing the microphone. The
implementation uses inference and tools rather than model training, keeps
local speech services under project ownership, and records the distinction
between local verification, prepared release, Pi deployment, and benchmark
evidence. The ElevenLabs adapter was provider-tested locally but is not part of
the active single-voice runtime. The final Qwen workstation acceptance passes
84 tests and the injection lane passes one clean fixture run. The Pi
browser-speaker path and vendored benchmark adapter pass local protocol checks,
including a small simulated 10-trial run. The Qwen Instruct audio profile is
provider-live over the restored VPN route: its matched run passes both ASR and
TTS, and the full injected voice/tool turn passes with zero errors. The next
decisive checks are automatic physical Pi speaker barge-in, operator-owned
live channel tests, and a clean captured v1.5 overlap
evaluation.

\begin{thebibliography}{9}

\bibitem{fdb}
Lin et al. ``FULL-DUPLEX-BENCH v1.5: Evaluating Overlap.''
Project source PDF: \textit{FULL-DUPLEX-BENCH V1.5 - EVALUATING OVERLAP.pdf},
AI 231 ME3 repository, 2026.

\bibitem{voiceagents}
``Building AI Voice Agents - Summary.'' Project source PDF, AI 231 ME3
repository, 2026.

\bibitem{whisper}
A. Radford et al. 2023. Robust speech recognition via large-scale weak
supervision. In \textit{Proceedings of ICML}.

\bibitem{silero}
Silero Team. Silero VAD. Open-source voice activity detection model and
runtime. The ME3 release pins the installed runtime version in
\texttt{requirements.txt}.

\bibitem{mcp}
Model Context Protocol. Official protocol specification and SDK. The ME3
release uses the Python MCP/FastMCP stack for owned task services.

\bibitem{pipecat}
Pipecat. Open-source framework for real-time voice and multimodal agents. The
ME3 release pins Pipecat 1.13.0 for the harness-facing adapter.

\bibitem{elevenlabs}
ElevenLabs. ``Realtime models'' and ``Stream speech.'' Official developer
documentation, 2026. \href{https://elevenlabs.io/docs/eleven-agents/customization/realtime-models}{ElevenLabs developer documentation}.

\bibitem{qwen2audio}
Qwen Team. ``Qwen2-Audio-7B-Instruct.'' Official model card, 2024.
\href{https://huggingface.co/Qwen/Qwen2-Audio-7B-Instruct}{Hugging Face model card}.

\bibitem{vllmtts}
vLLM Project. ``Qwen3-TTS online serving.'' Official vLLM-Omni recipe, 2026.
\href{https://github.com/vllm/vllm-omni/blob/main/recipes/Qwen/Qwen3-TTS.md}{vLLM-Omni recipe}.

\bibitem{resend}
Resend. ``Send Email.'' Official API reference, 2026.
\href{https://resend.com/docs/api-reference/emails/send-email}{Resend API reference}.

\end{thebibliography}

\end{document}

